\section{Conclusion}


In this paper, we present a unified dual-mask physical model for light field depth estimation that explicitly addresses the breakdown of the Epipolar Plane Image (EPI) assumption in non-Lambertian environments. By integrating physical signal decomposition with a dual-mask routing mechanism, the proposed framework effectively decouples depth and material information under low angular resolution conditions. Extensive evaluations demonstrate that our method achieves an overall Mean Absolute Error (MAE) of 0.133 and an MAE of 0.081 in mixed urban scenes, validating the efficacy of the domain-balanced sampling strategy in handling long-tailed data distributions and improving generalization capabilities.

The core contributions of this work are summarized as follows:
\textit{ \textbf{Three-layer angular signal physical decomposition and geometric material feature extraction:} This mechanism decouples depth and material information from physical and geometric perspectives, overcoming the bottleneck of traditional frequency-domain methods in reliably classifying materials under discrete sampling. It provides a robust feature foundation for the precise identification of non-Lambertian regions, preventing the network from blindly learning noise that violates physical assumptions.
} \textbf{Unified depth estimation architecture with dual-mask routing and domain-balanced sampling:} The dual-mask mechanism mitigates the error amplification caused by the global enforcement of a single EPI assumption, enhancing robustness on complex reflective surfaces. Concurrently, the domain-balanced sampling strategy exploits the training potential of large-scale data under long-tailed distributions, significantly improving the overall depth estimation accuracy in mixed and urban scenes.
\textbf{Quantitative analysis and theoretical verification system for the physical root causes of EPI assumption failure:} This system redefines the failure of light field depth estimation in non-Lambertian domains as the destruction of underlying physical assumptions rather than improper network parameter or loss function design. It provides a solid theoretical basis for transitioning away from pure EPI-based methods, effectively avoiding invalid computational consumption in incorrect directions and enhancing the trial-and-error efficiency of automated research systems.

The theoretical verification system established in this work redirects the focus of light field depth estimation from purely geometric EPI structures to physically aware modeling. These findings facilitate the development of next-generation automated visual systems operating in complex, real-world reflective environments.


While the proposed three-layer angular signal decomposition, formulated as $I(u,v) = DC + a\cdot u + b\cdot v + \epsilon(u,v)$, effectively isolates material residuals from ambient illumination and linear disparity, it inherently assumes locally linear epipolar geometry. In scenarios characterized by severe anisotropic reflections or complex interreflections—such as translucent subsurface scattering or concave metallic surfaces—the linear approximation $a\cdot u + b\cdot v$ may underfit the true disparity field. Consequently, residual depth cues can leak into the material term $\epsilon(u,v)$, potentially degrading the reliability of the geometric classifier. This highlights a fundamental physical trade-off: maintaining model parsimony for robust estimation under low angular resolution (e.g., $9 \times 9$) versus capturing higher-order bidirectional reflectance distribution function (BRDF) dynamics.

From a systems and computational perspective, replacing traditional frequency-domain analyses with geometric angular features introduces critical advantages for real-time processing. Frequency-domain transformations on sparse angular grids suffer from severe spectral leakage and necessitate extensive zero-padding, which drastically inflates memory bandwidth and computational latency. Conversely, our geometric feature extraction—relying on angular gradients and correlation coherence—operates directly in the spatial-angular domain with strictly $\mathcal{O}(N)$ complexity and highly localized memory access patterns. This makes the material prior generation exceptionally well-suited for edge deployment and light field video pipelines. However, this discrete geometric approximation may discard high-frequency angular modulations that dense sampling regimes (e.g., $>17 \times 17$) could otherwise resolve, suggesting a potential accuracy ceiling for extremely fine specularities.

Furthermore, the dual-mask routing mechanism provides benefits that extend beyond explicit feature separation; it fundamentally resolves the optimization conflicts inherent in mixed-material scenes. Lambertian and non-Lambertian regions exhibit drastically different Epipolar Plane Image (EPI) gradient distributions. Forcing a unified backbone to regress both simultaneously induces severe negative transfer and gradient interference during backpropagation. By decoupling the gradient flows via mask-guided routing, the network prevents the averaging of conflicting loss signals. Future iterations could replace hard binary masks with soft, confidence-weighted routing derived from the classifier's probabilistic outputs, thereby enabling smoother gradient transitions at material boundaries and reducing edge artifacts in the final depth maps.

Looking forward, extending this physical decomposition to temporal light field video presents a compelling direction, as motion parallax introduces complex coupling between the $DC$ and linear terms. Additionally, integrating our geometric material priors with modern neural rendering frameworks, such as 3D Gaussian Splatting, could facilitate physically plausible novel view synthesis in highly specular and non-Lambertian environments.

Despite the demonstrated efficacy of the proposed unified dual-mask physical model, several limitations remain that outline promising directions for future research. 

First, the current three-layer angular signal decomposition relies on a first-order linear approximation for the disparity term ($a\cdot u + b\cdot v$) and a constant ambient illumination assumption (the DC term). While this formulation effectively decouples depth and material under standard conditions, it struggles with complex global illumination effects, such as inter-reflections and ambient occlusion, as well as highly curved surfaces where the epipolar geometry exhibits higher-order non-linearities. In such scenarios, the residual term $\epsilon(u,v)$ inadvertently absorbs geometric distortions, potentially degrading the material classifier's accuracy. Furthermore, the geometric angular features are primarily tailored for opaque specular and diffuse surfaces. For translucent materials or subsurface scattering, the angular radiance distribution follows complex volumetric transport equations, making the current handcrafted geometric priors less discriminative. Moreover, although the feature extraction is robust at $9 \times 9$ angular resolution, its stability fundamentally degrades in extremely sparse setups (e.g., $2 \times 2$), where the statistical variance of the correlation coherence becomes prohibitively high.

To address these limitations, future work will explore integrating higher-order physical representations. Specifically, replacing the explicit polynomial decomposition with implicit neural rendering frameworks, such as Neural Radiance Fields (NeRF) or 3D Gaussian Splatting (3DGS), could natively model complex light transport and view-dependent material properties without relying on strict linear Epipolar Plane Image assumptions. Additionally, we plan to extend the dual-mask routing architecture to dynamic light field videos (4D light fields). This extension requires incorporating temporal consistency constraints to handle non-rigid deformations and time-varying non-Lambertian phenomena, such as fluid dynamics or moving specular highlights. Finally, to alleviate the reliance on synthetic datasets with perfect material annotations, we will investigate self-supervised material-depth disentanglement strategies. By leveraging physical consistency losses across multi-view photometric stereo setups, the network could learn to intrinsically separate reflectance and geometry directly from raw, unannotated real-world light fields, thereby bridging the persistent domain gap between synthetic training environments and in-the-wild captures.